\clearpage\fancyhead[L]{\small ANZEIGE.}\pdfbookmark[0]{Anzeigen. 1815 Juli 1.}{anzeige-1815-pfaff}\SourceLeaf{240}
\begin{center}\textbf{A N Z E I G E N.}\par\bigskip
\rule{.85\linewidth}{.4pt}\par
Göttingische gelehrte Anzeigen.\quad1815 Juli 1.\par
\rule{.85\linewidth}{.4pt}\end{center}
Der Königlichen Gesellschaft hat Herr Professor \textsc{Gauss} eine handschriftliche Abhandlung von Herrn Professor \textsc{Pfaff} in Halle mit folgendem Titel vorgelegt:
\begin{center}\textit{Methodus generalis, aequationes differentiarum partialium, nec non aequationes differentiales vulgares, utrasque primi ordinis, inter quotcunque variabiles, complete integrandi.}\end{center}
Die Theorie der partiellen Differentialgleichungen erster Ordnung verdankt \textsc{Lagrange} bekanntlich zwei wichtige Erweiterungen: ihre allgemeine Integration entweder dann, wenn sie bei beliebig vielen Variablen die partiellen Ableitungen nur linear enthalten, oder ohne Einschränkung ihrer Form, wenn es nur drei Variablen gibt. Über diese beiden Fälle war man bisher eigentlich noch nicht hinausgekommen, und das von \textsc{Lagrange} bei partiellen Differentialgleichungen in drei Variablen angewandte Verfahren würde bei einer größeren Anzahl besondere Schwierigkeiten bereiten. Wir müssen es daher als eine bemerkenswerte Bereicherung der Integralrechnung ansehen, dass es dem scharfsinnigen Verfasser der vorliegenden Abhandlung gelungen ist, die allgemeine Integration der partiellen Differentialgleichungen erster Ordnung für \textit{jede} Anzahl von Variablen zu finden. Er hat bei dieser Untersuchung einen eigenen Weg gewählt\SourceLeaf{241}%
 und sie mit einem anderen, nicht weniger interessanten Zweig der Integralrechnung verknüpft: der Theorie der gewöhnlichen Differentialgleichungen (erster Ordnung) zwischen mehr als zwei Variablen. Deren wahre Natur hat uns bekanntlich erst \textsc{Monge} erkennen gelehrt, obwohl dieser Geometer ihre Integration nur für die einfachsten Fälle vollendet hat. Für diese Art von Differentialgleichungen gibt Herr \textsc{Pfaff} die allgemeine Integration an; die Integration der partiellen Differentialgleichungen erscheint dann nur als ein besonderer Fall davon. Bei diesen Untersuchungen wird die allgemeine Integration der Differentialgleichungen jeder Ordnung zwischen zwei Variablen \textit{vorausgesetzt}. Das ist durchaus angemessen, ebenso wie man in den höheren Teilen der Mathematik die allgemeine Auflösung algebraischer Gleichungen postuliert.

\Needspace{12\baselineskip}
Wir glauben, den Freunden der höheren Mathematik einen willkommenen Dienst zu erweisen, wenn wir ihnen durch diese Anzeige jene schöne Erweiterung der Integralrechnung zugänglich machen. Ein Auszug aus der 144 Quartseiten umfassenden Abhandlung, in dem wir dem Verfasser Schritt für Schritt folgen und nichts Wesentliches übergehen wollten, würde freilich den uns zugestandenen Raum weit überschreiten. Wir wollen daher versuchen, uns allein an die \textit{Sache} zu halten und das Wesentliche in einer etwas veränderten Darstellung so herauszuarbeiten, dass Kenner es sich vollständig aneignen können.

Als Hauptoperation des gesamten Verfahrens ist die Reduktion eines Differentialausdrucks
\[
\Omega=p\,\mathrm dx+p'\mathrm dx'+p''\mathrm dx''+\text{usw.}+p^{(n-1)}\mathrm dx^{(n-1)}
\]
in dem jeder Koeffizient $p,p',p''$ usw. eine Funktion der Variablen $x,x',x''$ usw. ist, auf die Form
\[
\Omega=\lambda(q\,\mathrm dy+q'\mathrm dy'+q''\mathrm dy''+\text{usw.}+q^{(n-2)}\mathrm dy^{(n-2)});
\]
anzusehen, sodass $\lambda,y,y',y''$ usw. Funktionen von $x,x',x''$ usw. sind, dagegen $q,q',q''$ usw. Funktionen von $y,y',y''$ usw., und die Anzahl dieser letzteren Variablen um eins kleiner ist als die Anzahl der Variablen $x,x',x''$ usw.

Diese Transformation, die wir der Kürze halber mit (I) bezeichnen wollen, ist auf den Fall beschränkt, dass $n$ eine gerade Zahl ist. Wir werden weiter unten\SourceLeaf{242}%
 darlegen, wie sie auszuführen ist, und sie hier vorläufig voraussetzen, um den Überblick nicht zu unterbrechen.
\par
Wenn $n$ eine ungerade Zahl ist, wäre die Transformation I nur unter besonderen Bedingungen zwischen den Koeffizienten $p,p',p''.\,.\,.$ möglich; allgemein lässt sich in diesem Fall jedoch $\Omega$ auf die Form
\[
p^{*}\mathrm{d}x+\lambda(q\mathrm{d}y+q'\mathrm{d}y'+q''\mathrm{d}y''+\mathrm{usw.}
+q^{(n-3)}\mathrm{d}y^{(n-3)})
\]
bringen. Man betrachte dazu zunächst $x$ in $\Omega$ als konstant und transformiere unter dieser Voraussetzung nach (I) den Ausdruck
\[
p'\mathrm{d}x'+p''\mathrm{d}x''+p'''\mathrm{d}x'''+\mathrm{usw.}
+p^{(n-1)}\mathrm{d}x^{(n-1)}
\]
in dem die Anzahl der Variablen $x',x'',x'''.\,.\,.$ nun gerade ist, in
\[
\lambda(q\mathrm{d}y+q'\mathrm{d}y'+q''\mathrm{d}y''+\mathrm{usw.}
+q^{(n-3)}\mathrm{d}y^{(n-3)})=\lambda\Omega'
\]
\Needspace{13\baselineskip}
Hier sind also $q,q',q''$ usw. Funktionen von $y,y',y''$ usw.; diese wiederum sind, ebenso wie $\lambda$, Funktionen von $x,x',x'',x'''.\,.\,.\,.$. Von diesen Größen muss jedoch die erste, $x$, als konstant behandelt werden, um aus der Entwicklung von $\lambda\Omega'$ den Ausdruck
\[
p'\mathrm{d}x'+p''\mathrm{d}x''+p'''\mathrm{d}x'''+\mathrm{usw.}
\]
zu erhalten. Das zusätzliche Glied, das entsteht, wenn bei jener Entwicklung auch $x$ als variabel betrachtet wird, ist
\[
=(q\cdot\frac{\mathrm{d}y}{\mathrm{d}x}
+q'\cdot\frac{\mathrm{d}y'}{\mathrm{d}x}
+q''\cdot\frac{\mathrm{d}y''}{\mathrm{d}x}+\mathrm{usw.})\lambda\mathrm{d}x
\]
Um die obige Form zu erhalten, braucht man daher nur
\[
p^{*}=p-\lambda(q\cdot\frac{\mathrm{d}y}{\mathrm{d}x}
+q'\cdot\frac{\mathrm{d}y'}{\mathrm{d}x}
+q''\cdot\frac{\mathrm{d}y''}{\mathrm{d}x}+\mathrm{usw.})
\]
zu setzen. Diese Transformation von $\Omega$ in $p^{*}\mathrm{d}x+\lambda\Omega'$, die auf ungerade Werte von $n$ beschränkt ist, wollen wir mit II bezeichnen. Offenbar kann dieselbe Reduktion erneut auf $\Omega'$ angewandt und
\[
\begin{aligned}
&\Omega'=q^{*}\mathrm{d}y+\lambda'(r\mathrm{d}z+r'\mathrm{d}z'+r''\mathrm{d}z''+\mathrm{usw.}\\&\qquad+r^{(n-5)}\mathrm{d}z^{(n-5)})\\
&=q^{*}\mathrm{d}y+\lambda'\Omega''
\end{aligned}
\]
\SourceLeaf{243}
gesetzt werden, dann wiederum $\Omega''=r^{*}\mathrm{d}z+\lambda''\Omega'''$, bis man schließlich zu einem Ausdruck gelangt, der nur eine Variable enthält. Damit ist also $\Omega$ auf die Form
\[
p^{*}\mathrm{d}x+\lambda q^{*}\mathrm{d}y+\lambda\lambda'r^{*}\mathrm{d}z+\mathrm{usw.}
\]
oder auf die Form
\[
P\mathrm{d}x+Q\mathrm{d}y+R\mathrm{d}z+\mathrm{usw.}
\]
gebracht. Hier gilt für die Anzahl der Variablen $x,y,z$ usw. $=\frac12(n+1)$, und sämtliche $n$ Größen $y,z$ usw., $P,Q,R$ usw. sind Funktionen von $x,x',x''$ usw. Dieses Reduktionsverfahren sei mit III bezeichnet.

\Needspace{19\baselineskip}
Wendet man dieses Verfahren III im Fall einer ursprünglich gegebenen geraden Anzahl von Variablen auf den durch Reduktion I erhaltenen Ausdruck
\[
q\mathrm{d}y+q'\mathrm{d}y'+q''\mathrm{d}y''+\mathrm{usw.}
\]
an, so gelangt dadurch
\[
\begin{aligned}&\Omega=p\mathrm{d}x+p'\mathrm{d}x'+p''\mathrm{d}x''+\mathrm{usw.}
+p^{(n-1)}\mathrm{d}x^{(n-1)}\\&\qquad\text{in die Form}\end{aligned}
\]
\[
Q\mathrm{d}y+R\mathrm{d}z+\mathrm{usw.}
\]
sodass die Anzahl der Variablen $y,z$ usw. $=\frac12 n$ gilt und alle $n$ Größen $Q,R$ usw., $y,z$ usw. Funktionen von $x,x',x''$ usw. werden. Diese Reduktion sei mit IV bezeichnet.

Diese allgemeine Transformierbarkeit der Differentialausdrücke nach III und IV ist ein ebenso neuer wie bemerkenswerter Lehrsatz. Er wird in der Abhandlung von Herrn \textsc{Pfaff} zwar nicht ausdrücklich ausgesprochen, lässt sich aber leicht aus den dortigen Untersuchungen folgern.

Daraus lassen sich nun die Lösungen der am Anfang dieser Anzeige genannten Aufgaben leicht ableiten.

\Needspace{17\baselineskip}
\medskip
1) Um die Differentialgleichung
\[
0=p\mathrm{d}x+p'\mathrm{d}x'+p''\mathrm{d}x''+\mathrm{usw.}
\]
oder
\[
0=\Omega
\]
\SourceLeaf{244}
zu integrieren, in der $p,p',p''$ usw. gegebene Funktionen der $n$ Variablen $x,x',x''$ usw. sind, setzt man bei geradem $n$ nach IV
\[
\Omega=Q\mathrm{d}y+R\mathrm{d}z+S\mathrm{d}u+\mathrm{usw.}
\]
wobei $Q,y.R,z,S,u$ usw. zusammen $n$ gegebene Funktionen von $x,x',x''$ sind. Die Differentialgleichung
\[
0=Q\mathrm{d}y+R\mathrm{d}z+S\mathrm{d}u+\mathrm{usw.}
\]
ist also der vorgegebenen gleichwertig, und ihre allgemeinste Integration ist im folgenden System von $\frac12 n$ Gleichungen enthalten:
\[
\begin{aligned}&0=\varphi(y,z,u\ \mathrm{usw.}),\\&
\frac1Q\cdot\frac{\mathrm{d}\varphi(y,z,u.\,.\,.)}{\mathrm{d}y}
\\&=\frac1R\cdot\frac{\mathrm{d}\varphi(y,z,u.\,.\,.)}{\mathrm{d}z}
\\&=\frac1S\cdot\frac{\mathrm{d}\varphi(y,z,u.\,.\,.)}{\mathrm{d}u}=\mathrm{usw.}\end{aligned}
\]
Dabei bezeichnet $\varphi$ eine beliebige Funktion, und die Ableitungen sind selbstverständlich \textit{partielle}. Soweit durch die Gleichung
$0=\varphi(y,z,u\ \mathrm{usw.})$
die Größe $y$ durch die übrigen bestimmt werden kann, lässt sich die Lösung auch durch folgende Gleichungen darstellen:
\[
\begin{gathered}
y=\psi(z,u.\,.\,.\,.)\\
\frac{\mathrm{d}\psi(z,u.\,.\,.\,.)}{\mathrm{d}z}=-\frac RQ\\
\frac{\mathrm{d}\psi(z,u.\,.\,.\,.)}{\mathrm{d}u}=-\frac SQ\\
\text{u.\,s.\,f.}
\end{gathered}
\]
Genau genommen wäre diese Lösung jedoch weniger allgemein, da die beliebige Funktion $\varphi(y,z,u.\,.\,.\,.)$ auch solche umfasst, in denen $y$ nicht vorkommt.

\Needspace{17\baselineskip}
\medskip
2) Zur Integration derselben Differentialgleichung bei ungeradem $n$ setzt man $\Omega$ nach III in folgende Form:
\[
\Omega=P\mathrm{d}x+Q\mathrm{d}y+R\mathrm{d}z+\mathrm{usw.}
\]
Hier sind $P,Q,y,R,z$ usw. zusammen $n$ gegebene Funktionen von $x,x',x''$ usw. Die allgemeinste Integration der Differentialgleichung $\Omega=0$ beruht dann auf folgendem System von $\frac12(n+1)$ Gleichungen:
\[
\begin{aligned}&0=\varphi(x,y,z.\,.\,.),\\&
\frac1P\cdot\frac{\mathrm{d}\varphi(x,y,z.\,.\,.)}{\mathrm{d}x}
=\frac1Q\cdot\frac{\mathrm{d}\varphi(x,y,z.\,.\,.)}{\mathrm{d}y}
\\&=\frac1R\cdot\frac{\mathrm{d}\varphi(x,y,z.\,.\,.)}{\mathrm{d}z}=\mathrm{usw.}\end{aligned}
\]
\par\SourceLeaf{245}
3) Die allgemeine Integration einer gegebenen partiellen Differentialgleichung erster Ordnung, also einer endlichen Gleichung zwischen den partiellen Ableitungen
\[
\frac{\mathrm{d}x}{\mathrm{d}x'}=p',\qquad
\frac{\mathrm{d}x}{\mathrm{d}x''}=p'',\qquad
\frac{\mathrm{d}x}{\mathrm{d}x'''}=p''',\quad\text{u.\,s.\,f.}
\]
und $x,x',x'',x'''$ usw. (wobei $x$ eine noch zu bestimmende Funktion der $m$ Variablen $x',x'',x'''$ usw. bezeichnet), ist nichts anderes als die allgemeine Integration der gewöhnlichen Differentialgleichung
\[
0=-\mathrm{d}x+p'\mathrm{d}x'+p''\mathrm{d}x''+p'''\mathrm{d}x'''+\text{usw.}
\]
Da aufgrund jener endlichen Gleichung eine der Größen $p',p'',p'''$ usw., etwa $p'$, als Funktion der übrigen Größen $p'',p'''$ usw. und von $x,x',x'',x'''$ usw. dargestellt werden kann, ist die eben angegebene Differentialgleichung als eine zwischen den $2m$ Variablen $x,x',x''$ usw., $p'',p'''$ usw. zu betrachten, in der die Differentiale $\mathrm{d}p'',\mathrm{d}p'''$ usw. den Koeffizienten $0$ haben. Um die Integration auszuführen, bringt man daher den Differentialausdruck
\[
-\mathrm{d}x+p'\mathrm{d}x'+p''\mathrm{d}x''+p'''\mathrm{d}x'''+\text{usw.}
\]
auf die Form
\[
Q\mathrm{d}y+R\mathrm{d}z+S\mathrm{d}u+\text{usw.}
\]
wobei die $2m$ Größen $Q,y,R,z,S.\ u$ usw. bekannte Funktionen von $x,x',x''\, .\, .\, .\, .\, p'',p'''\, .\, .\, .\, .$ sind. Die Integration ist dann in demselben Gleichungssystem wie oben (1) enthalten. Denkt man sich daraus $p'',p'''$ usw. eliminiert, so bleibt eine endliche Gleichung zwischen $x,x',x''.x'''$ usw. zurück. Die \emph{wirkliche} Elimination kann freilich nur ausgeführt werden, soweit für $\varphi$ \emph{bestimmte} Funktionen angenommen werden. Dieser Umstand beruht jedoch auf der Natur des Problems und nicht auf der Unvollkommenheit der Analyse, die, solange sie beim Allgemeinen bleibt, die Lösung nur in jener Form geben kann.

Im Übrigen sieht man unmittelbar, dass wir auf ähnliche Weise auch \emph{mehrere} gleichzeitig bestehende partielle Differentialgleichungen integrieren können.


